\section{Síntesis: características que determinan el comportamiento}
\label{sec:sintesis}

Esta sección recoge, a partir del análisis anterior, las características de diseño que
condicionan los resultados de ambos métodos. Se distinguen las que son comunes, y por tanto
explican el nivel de calidad que alcanzan los dos, de las que los diferencian.

\subsection{Características comunes}

\begin{enumerate}
\item \textbf{Rutas siempre factibles.} Las rutas nunca violan capacidad ni ventanas de
tiempo. La única infactibilidad que admite la búsqueda son los clientes sin asignar, que
tienen un peso fijo de $10^6$ en $f$; no hay otras penalizaciones que calibrar.
\item \textbf{Objetivo jerárquico escalarizado.} Los pesos de $10^6$ por cliente sin asignar y
de $10^4$ por vehículo hacen que colocar un cliente domine sobre cualquier reducción de flota,
y esta sobre cualquier mejora de distancia, en la aceptación y en la detección de mejoras.
\item \textbf{Reducción de flota en una fase propia.} La reparación no abre rutas, de modo que
la flota es fija en cada etapa. La fase~1 elimina rutas una a una (retira la más corta y
busca recolocar sus clientes con el ALNS completo) y se detiene al alcanzar la cota inferior
de capacidad, al estancarse con $5$ o más clientes pendientes durante $2000$ iteraciones o al
agotar $T_{\max}$ iteraciones. La fase~2 optimiza la distancia durante otras $T_{\max}$
iteraciones con la flota resultante. El número de vehículos obtenido depende, por tanto, de
la eficacia con que cada método selecciona operadores \emph{durante la fase~1}.
\item \textbf{Vecindarios de tamaño acotado.} Se extrae entre el 10\,\% y el 40\,\% de los
clientes, con topes de $30$ y $60$: a partir de $300$ clientes el vecindario tiene el mismo
tamaño absoluto, lo que mantiene estable el costo por iteración al crecer la instancia.
\item \textbf{Catálogo complementario.} Hay operadores de diversificación (extracción
aleatoria), de intensificación (extracción de peores, inserción voraz), de estructura (Shaw,
grupos, rutas), específicos de las ventanas de tiempo y con anticipación (arrepentimiento 2,
3, 4 y sobre todas las rutas). El ruido en los costos de inserción, activo en la mitad de las
iteraciones, diversifica cualquiera de las reparaciones.
\item \textbf{Recocido con reinicios.} La temperatura proporcional a la distancia inicial,
el enfriamiento geométrico y el retorno a la mejor solución con recalentamiento a
$0.5\,T_0$ tras $1500$ iteraciones sin progreso evitan que la búsqueda se congele en los
periodos de estancamiento. La temperatura vuelve a $T_0$ al comenzar cada etapa de la fase~1
y al comenzar la fase~2.
\item \textbf{Eficiencia por iteración.} La evaluación de inserciones en tiempo constante, las
cachés por ruta y las estructuras incrementales de los operadores de destrucción reducen el
costo de la reparación en un factor del orden del número de rutas. Al ser optimizaciones
exactas, no modifican la trayectoria de búsqueda: lo que aportan es que un presupuesto de
hasta $50\,000$ iteraciones por corrida sea alcanzable en las instancias de mayor tamaño.
\end{enumerate}

\subsection{Diferencias entre los dos métodos}

Como ambos métodos comparten todo lo anterior y parten de la misma solución inicial con la
misma semilla, cualquier diferencia de resultados se origina en los aspectos de la
Tabla~\ref{tab:comparacion}.

\begin{table}[H]
\centering
\small
\caption{Comparación de los mecanismos de selección y aprendizaje.}
\label{tab:comparacion}
\begin{tabularx}{\textwidth}{@{}l L L@{}}
\toprule
Aspecto & ALNS clásico & ALNS con Q-learning \\
\midrule
Unidad de decisión & Destrucción y reparación por separado ($6+5$ pesos) & Par conjunto ($30$ acciones) \\
Contexto & Ninguno & Estado binario: si la iteración anterior mejoró \\
Política & Sorteo proporcional a los pesos & $\varepsilon$-voraz, $\varepsilon\to 0.01$ \\
Señal de crédito & Categoría del resultado: $33$, $13$, $9$, $0$ & Magnitud de la mejora relativa en su propio criterio, logarítmica, ponderada global/local y creciente en el tiempo \\
Aceptación sin mejora & Se premia ($9$) & Se penaliza igual que un rechazo \\
Iteración sin éxito & Puntaje $0$ & Recompensa negativa $-\tau\,\Omega$: la mejora media reciente del mejor par alternativo \\
Frecuencia de actualización & Cada $100$ iteraciones & Cada iteración \\
Memoria & Media móvil, $\lambda=0.9$ por segmento & Media móvil, $\alpha=0.5$ por visita \\
Horizonte & Inmediato & Descontado, $\gamma=0.7$ \\
Exploración & Permanente y proporcional ($w_{\min}=0.01$) & $200$ iteraciones uniformes al inicio de cada fase y después residual \\
Cambio de fase & Sin reacción: los pesos continúan & Olvida las mejoras recientes y reinicia la exploración; conserva la tabla $Q$ \\
Reparto de uso & Gradual entre todos los operadores & Concentrado en el par mejor valorado de cada estado \\
\bottomrule
\end{tabularx}
\end{table}

De la tabla se desprenden cuatro diferencias de comportamiento que conviene tener presentes al
interpretar los resultados experimentales:

\begin{itemize}
\item \textbf{Sinergias entre operadores.} Solo el Q-learning puede aprender que un operador
de destrucción rinde con un operador de reparación concreto: cada una de las seis
destrucciones se valora por separado con cada una de las cinco reparaciones.
\item \textbf{Qué se entiende por éxito.} El método clásico premia la aceptación, de modo que
a temperatura alta refuerza a los operadores que generan soluciones aceptables aunque no
mejoren. El Q-learning solo premia la mejora estricta y pondera su magnitud relativa, y
penaliza cada fallo en proporción a lo que están rindiendo los demás pares.
\item \textbf{Velocidad y concentración de la adaptación.} El método clásico modifica su
distribución lentamente y mantiene en uso todo el catálogo; el Q-learning reacciona en pocas
visitas y explota casi en exclusiva el par preferido de cada estado, cambiándolo cuando las
penalizaciones lo desplazan. El primero es más estable y diverso; el segundo, más reactivo y
más intensivo.
\item \textbf{Adaptación al cambio de objetivo.} Al pasar de eliminar rutas a reducir
distancia, el método clásico arrastra los pesos de la fase~1 y los corrige al ritmo de su
media móvil; el Q-learning vuelve a explorar todos los pares de forma explícita antes de
explotar. La fase~1 es, además, la que decide el número de vehículos, y en ella ambos
mecanismos disponen de un número de iteraciones variable (a lo sumo $T_{\max}$) para aprender.
\end{itemize}

Estas afirmaciones describen el mecanismo, no su efecto medido. Para vincularlas con los
resultados se recomienda acompañar las tablas comparativas con la frecuencia de uso de cada
operador (o par) a lo largo de la búsqueda en ambos métodos, separada por fases, y con la
duración de la fase~1.
